\section{Open Questions}

The replication verified the paper's headline claims (C3--C6) for local depolarizing noise on the
specific TIsing $n{=}4$ system, but leaves five substantive follow-ups open. Each is stated with
its basis and a concrete next-step recipe implementable within the existing harness.

\subsection*{Q1. Correlated / coherent noise beyond independent depolarizing}
\textbf{Basis.} The paper's Eq.~(6) noise model assumes independent single-gate Kraus channels; our
replication inherits this assumption. Real superconducting hardware exhibits ZZ crosstalk between
neighboring qubits and coherent unitary errors from imperfect pulse calibration, both of which can
violate the linear-accumulation ansatz. This gap was not addressed by the paper or our replication.
\\
\textbf{Next steps.} In the existing harness, replace
\texttt{add\_all\_qubit\_quantum\_error(depolarizing\_error(p,k))} with
(a) a \texttt{QuantumChannel} built from two-qubit correlated Pauli errors on nearest-neighbor pairs
at strength $p$, and
(b) coherent $R_x/R_z$ overrotations at angle $\sqrt{p}$.
Rerun the $n{=}4$, $d{=}2,3$ sweep and fit $\Delta E(p)$ to $p$, $p\log p$, and $p^\alpha$. Report
deviation from linearity and whether per-gate-slope intensiveness (C6) is preserved.

\subsection*{Q2. Integration with modern error mitigation (ZNE, PEC, CDR)}
\textbf{Basis.} The paper predates the widespread adoption of ZNE/PEC (Kandala et al.\ 2019, Temme
et al.\ 2017, Czarnik et al.\ 2021) and does not attempt any mitigation. Whether the observed linear
$\Delta E(p)$ can be canceled by Richardson extrapolation is a distinct empirical question with direct
hardware relevance.
\\
\textbf{Next steps.} Using the existing $d{=}2$ sweep at $p \in \{10^{-4}, 3\!\cdot\!10^{-4},
10^{-3}, 3\!\cdot\!10^{-3}, 10^{-2}\}$: fit $E(p)$ to a low-order polynomial in $p$ and extrapolate
to $p{=}0$ via Richardson. Compare extrapolated $E(0)$ against our measured $E(0){=}-5.0466$ and
against $E_0{=}-5.2263$. Report residual error vs mitigation order. Then rerun with folded-CNOT ZNE
at noise scales $\lambda \in \{1, 3, 5\}$ at fixed $p{=}10^{-3}$ (\texttt{mitiq}-compatible).

\subsection*{Q3. Connection to noise-induced barren plateaus (NIBP)}
\textbf{Basis.} Wang et al.\ (Nat.\ Commun.\ 2021) showed that local Pauli noise flattens the VQE
cost landscape exponentially in circuit depth. The paper we replicated does not address gradient
variance, and neither does our fixed-$\theta^\star$ protocol. Our observed linear $\Delta E(p)$ is
measured at a fixed optimal noiseless point, not at a re-optimized noisy minimum.
\\
\textbf{Next steps.} Modify \texttt{vqe\_noisy.py} to re-run COBYLA \emph{inside} the noise loop
(cost function $=$ Tr$(\rho(\theta) H)$ at each $p$). Compute (a) the variance
$\text{Var}[\partial E/\partial \theta_i]$ over 20 random $\theta$ samples at each $p$, and (b) the
noisy re-optimized energy $E_{\min}(p)$. Test whether Var falls exponentially in $d$ for
$d \in \{2,3,4,5\}$ at $p{=}10^{-3}$, matching the NIBP prediction Var $\sim e^{-c d}$.

\subsection*{Q4. Device-specific composite noise models}
\textbf{Basis.} The paper's Sec.~IV composite IBM model (their C9, not tested here) is the only
device-specific comparison in the original work. Newer NISQ devices (2024--2026 vintage) have $T_1/T_2$
imbalance and mid-circuit measurement noise that the depolarizing model does not capture. Whether the
paper's scaling laws survive under hardware realism is an open, hardware-relevant question.
\\
\textbf{Next steps.} Build three \texttt{NoiseModel} instances in Qiskit Aer using calibration data
pulled from \texttt{IBMProvider().get\_backend('ibm\_...').properties()} (public FakeBackend snapshots
suffice), a synthetic Rigetti-style crosstalk model, and an ion-trap $T_1$-limited model with
$T_1{=}10\,\text{s}, T_2{=}1\,\text{s}$. Re-run the $n{=}4$, $d{=}2,3$ sweep against each. Report
per-device linear slopes and whether the depth-accumulation ratio remains close to the gate-count
ratio.

\subsection*{Q5. Noise-aware ansatz redesign}
\textbf{Basis.} The paper adopts a generic hardware-efficient ansatz and does not attempt noise-aware
ansatz design. Recent work (McClean et al.\ 2018, Grimsley et al.\ 2019 ADAPT-VQE, Bharti et al.\ 2022
review) argues that problem-adapted ansatze can achieve comparable expressivity with fewer gates,
directly reducing the linear-slope coefficient. This is actionable and testable in our harness.
\\
\textbf{Next steps.} Implement (a) a hardware-efficient ansatz restricted to the $\mathbb{Z}_2$-even
subspace (parity-preserving CNOT + $R_{zz}$ + $R_x$ layers only, $\sim 15$ gates at $n{=}4, d{=}2$),
and (b) an ADAPT-VQE built from a Pauli pool of $\{Z_iZ_j, X_i\}$ operators grown greedily until
noiseless $E/E_0 \ge 0.98$. Measure the linear $\Delta E(p)$ slope for each; report
slope-per-gate vs slope-per-VQE-precision tradeoff. A slope reduction of $>20\%$ at fixed noiseless
fidelity would validate noise-aware design as a lever.
